\chapter{Deep Dive: JWT at Scale --- Where to Verify, Fat Tokens, and Asymmetric Keys}
\label{ch:jwt-scale}

Chapter~5 taught the project's security design as it is: HS256 tokens,
verified once at the gateway, identity forwarded as headers. That design
is correct for a system with one gateway and five services. This chapter
is the first of the book's \emph{deep dives}: it takes one design already
in the project and pushes it until it breaks, the way an interviewer or a
tenfold traffic increase would. The questions come from a widely shared
Vietnamese article on JWT in
microservices;\footnote{``Lột trần sự thật về JWT --- Phần 2: Trận chiến
Microservices và nỗi ám ảnh Token béo phì'', viblo.asia, 2026.} the
answers come from reading the project's code.

Three questions organize the chapter:

\begin{enumerate}
  \item \textbf{Where should a token be verified} --- at the gateway
  once, or in every service?
  \item \textbf{What happens when the token gets fat} --- and what is the
  \emph{phantom token} pattern that keeps it thin?
  \item \textbf{What does a signature actually guarantee}, and why does
  the choice between HS256 and RS256 decide who can be trusted with what?
\end{enumerate}

\section{Where to verify: gateway or every service?}
\label{sec:where-to-verify}

\subsection{Model 1: centralized verification (what the project does)}

Re-read the two halves of the trusted-edge pattern. At the edge, in the
gateway's \code{JwtAuthenticationFilter}:

\begin{lstlisting}[language=Java]
Claims claims = Jwts.parser()
        .verifyWith(signingKey)                    (*@\co{1}@*)
        .build()
        .parseSignedClaims(token)
        .getPayload();

String userId = claims.getSubject();
String role = claims.get(JwtConstants.CLAIM_ROLE, String.class);

ServerHttpRequest mutatedRequest = exchange.getRequest().mutate()
        .header(JwtConstants.HEADER_USER_ID, userId)   (*@\co{2}@*)
        .header(JwtConstants.HEADER_USER_ROLE, role)
        .build();
\end{lstlisting}

And behind it, in every service's \code{GatewayHeaderAuthFilter}
(course-service's copy shown):

\begin{lstlisting}[language=Java]
String userId = request.getHeader("X-User-Id");
String userRole = request.getHeader("X-User-Role");

if (userId != null && userRole != null) {               (*@\co{3}@*)
    var authorities = List.of(
            new SimpleGrantedAuthority("ROLE_" + userRole));
    var authentication = new UsernamePasswordAuthenticationToken(
            Long.valueOf(userId), null, authorities);
    SecurityContextHolder.getContext().setAuthentication(authentication);
}
\end{lstlisting}

\begin{description}
  \item[\co{1}] Cryptography happens exactly once, here. The key is the
  same \code{app.jwt.secret} that auth-service signs with.
  \item[\co{2}] Identity leaves the gateway as two plain HTTP headers.
  Nothing downstream can tell whether the gateway wrote them or someone
  else did.
  \item[\co{3}] The service performs \emph{no} check. Two headers
  present means authenticated. This is the whole contract.
\end{description}

The advantages are real: one place holds the secret, services stay free of
JWT libraries and key handling, and the per-request CPU cost of signature
checking is paid once. Most Spring Cloud tutorials stop here.

\subsection{The ``trusted internal network'' fallacy}

The pattern rests on one invariant --- \emph{every} path to a service
passes through the gateway. In the project's Compose file that invariant
is already false. \code{docker-compose.yml} publishes auth-service on
\code{8081} and course-service on \code{8082} directly to the host, so
that you can debug them. Anyone on the host can therefore skip the
gateway:

\begin{handson}[forge an identity in one line]
With the stack running, ask course-service for a teacher's courses without
ever presenting a token:
\begin{lstlisting}[language=bash]
$ curl -s -H "X-User-Id: 1" -H "X-User-Role: TEACHER" \
       http://localhost:8082/courses | head -c 300
\end{lstlisting}
It answers. Now repeat through the gateway on \code{8080} with the same
two headers and no \code{Authorization}:
\begin{lstlisting}[language=bash]
$ curl -s -H "X-User-Id: 1" -H "X-User-Role: TEACHER" \
       http://localhost:8080/api/courses
{"error":"Unauthorized","message":"Missing or invalid Authorization header"}
\end{lstlisting}
The difference between the two outputs is the entire security model.
\end{handson}

In Kubernetes the published ports go away (Chapter~18), but the flat pod
network remains: any pod in the namespace can reach
\code{course-service:8082}. The fallacy is assuming that ``inside'' means
``trusted''. One compromised container --- a vulnerable image, a leaked
kubeconfig, a supply-chain dependency --- and the attacker forges
\code{X-User-Id: 1} to every service at once.

\begin{pitfall}[identity headers survive public paths]
Symptom: nothing visible, which is the problem. Look at the top of
\code{JwtAuthenticationFilter.filter}: a request to a public path
(\code{/api/auth/refresh}, \code{/actuator/**}) is forwarded
\emph{unmodified} --- the gateway neither verifies a token nor touches
headers. A client can therefore send \code{X-User-Id: 1} to a public path
and the header arrives at the service intact. Today no public endpoint
reads it, so nothing is exploitable; but the invariant ``only the gateway
writes these headers'' is not enforced, it is merely true by luck. The fix
is one default filter in \code{api-gateway.yml} that strips both headers
from every \emph{incoming} request before the JWT filter sets them:
\begin{lstlisting}[language=yaml]
spring:
  cloud:
    gateway:
      default-filters:
        - RemoveRequestHeader=X-User-Id
        - RemoveRequestHeader=X-User-Role
\end{lstlisting}
Now the headers can only originate at the gateway. This closes the
external door; it does nothing for the in-cluster one.
\end{pitfall}

\subsection{Model 2: zero trust --- every service verifies}

The alternative removes the trust: every service parses and verifies the
JWT itself, and ignores any identity header. A forged header is now
worthless because the service never reads it; a forged token fails
signature verification. The gateway may still verify --- rejecting bad
tokens early is cheap --- but nothing downstream \emph{depends} on it.

The cost appears immediately: every verifier needs the key. With HS256
that means copying the \emph{signing} secret to every service, which is
worse than the problem it solves (Section~\ref{sec:hs-vs-rs}). Zero trust
at the application layer is only reasonable with asymmetric keys, which
is why the two halves of this chapter belong together.

\begin{decision}[gateway-verified headers, for now]
The project keeps Model~1. It has exactly two key holders, no
third-party services, and every service is deployed by the same
pipeline --- the blast radius of the trusted-edge assumption is the
single namespace. Three signals would flip the decision: a service
written by another team, a service reachable by any caller other than
the gateway (a Kafka consumer that also exposes HTTP, a batch job), or
a compliance requirement that identity be verifiable at the point of
use. At that moment the move is RS256 plus per-service verification,
not more headers. Production systems usually take a third path: keep the
headers, and make the \emph{transport} prove the sender with mutual TLS
from a service mesh (Istio, Linkerd). The gateway then becomes a
cryptographically identified peer instead of an unverifiable
neighbor.
\end{decision}

\section{What a signature guarantees}

A JWT is \code{base64url(header).base64url(payload).signature}. The first
two parts are \emph{encoded}, not encrypted: paste any project token into
a decoder and \code{sub}, \code{email}, \code{role} are readable. The
signature adds exactly two properties:

\begin{itemize}
  \item \textbf{Integrity.} Change one byte of the payload
  (\code{"role":"STUDENT"} to \code{"TEACHER"}) and verification fails.
  \item \textbf{Authenticity.} The token was produced by a holder of the
  signing key --- and by nobody else.
\end{itemize}

It adds no confidentiality. Both HS256 and RS256 compute the signature
over the same string, \code{base64url(header) + "." +
base64url(payload)}. They differ only in \emph{which key creates the
signature and which key checks it}. That single difference decides the
whole trust topology of a system.

\section{HS256 versus RS256}
\label{sec:hs-vs-rs}

\subsection{HS256: one key, two roles}

HS256 is HMAC-SHA256: a hash with a key mixed in. You do not need the
construction (two nested SHA-256 passes with the key XORed against fixed
pads); you need one property. \emph{The same key both creates and checks
the signature.} Verifying means recomputing the HMAC with your copy of
the key and comparing bytes. Without the key you cannot compute it, so you
cannot forge. With the key you can compute it, so you \emph{can} forge.
Verifier and signer are indistinguishable.

The project's two copies of that key:

\begin{lstlisting}[language=Java]
// auth-service/.../service/JwtService.java
this.signingKey = Keys.hmacShaKeyFor(secret.getBytes(StandardCharsets.UTF_8));

// api-gateway/.../filter/JwtAuthenticationFilter.java
this.signingKey = Keys.hmacShaKeyFor(jwtSecret.getBytes(StandardCharsets.UTF_8));
\end{lstlisting}

Both read \code{app.jwt.secret}; in the cluster both mount the same
\code{jwt} Secret (Chapter~16). The gateway only ever needs to
\emph{verify}, yet it holds a key that can \emph{mint} a token for any
user with any role. Compromise the gateway pod and you own every identity
in the system. Extend Model~2 naively --- copy the secret to five
services --- and you have six such places.

\subsection{RS256: two keys, one role each}

RS256 is an RSA signature over a SHA-256 digest. RSA rests on an
asymmetry in arithmetic: multiplying two large primes is instant;
recovering them from the product is infeasible.

Key generation, once:

\begin{enumerate}
  \item Pick two large primes $p$, $q$ (about 1024 bits each for a
  2048-bit key). Compute $n = p \times q$.
  \item Pick a public exponent $e$; in practice always $65537$. The pair
  $(n, e)$ is the \textbf{public key}.
  \item Compute $d$ such that $e \times d \equiv 1 \pmod{(p-1)(q-1)}$.
  The pair $(n, d)$ is the \textbf{private key}. Computing $d$ requires
  $p$ and $q$; knowing only $n$ does not give them.
\end{enumerate}

Signing and verifying:

\begin{lstlisting}
h = SHA256(header "." payload), padded (PKCS#1 v1.5)
sign:    signature = h^d  mod n      -- needs d: only auth-service
verify:  signature^e mod n == h ?    -- needs e, n: anyone
\end{lstlisting}

Because $e \times d \equiv 1$, raising to $d$ and then to $e$ returns
$h$. Anyone with the public key can check a signature; nobody without $d$
can produce one. That is the asymmetry HS256 lacks, and it changes who
may hold what:

\begin{figure}[htb]
\centering
\begin{tikzpicture}[node distance=6mm and 15mm]
  \node[boxdark] (auth) {auth-service\\private key $(n,d)$};
  \node[boxfill, right=of auth] (jwks) {JWKS endpoint\\public key $(n,e)$, \code{kid}};
  \node[box, right=of jwks] (gw) {api-gateway\\verify only};
  \node[box, above=of gw] (dict) {dictionary-service\\verify only};
  \node[box, below=of gw] (course) {course-service\\verify only};
  \draw[flow] (auth) -- node[lbl,above]{publishes} (jwks);
  \draw[flowdash] (jwks.east) -- node[lbl,below,pos=0.55]{fetch, cache} (gw.west);
  \draw[flowdash] (jwks.east) -- (dict.west);
  \draw[flowdash] (jwks.east) -- (course.west);
\end{tikzpicture}
\caption{RS256 key topology: one signer, any number of verifiers. Leaking
a verifier leaks nothing.}
\label{fig:rs256-topology}
\end{figure}

\begin{itemize}
  \item Auth-service holds the private key and is the \emph{only}
  place a token can be issued.
  \item Gateway and every service hold the public key. Leaking it is
  harmless --- it is public by definition. You could print it in this
  book.
  \item Key distribution stops being a secret-management problem: publish
  the public key at a well-known URL (a \emph{JWKS}, JSON Web Key Set),
  and any verifier fetches it at startup and caches it.
\end{itemize}

You already consume this topology. The Google login of Chapter~5 returns
an \code{id\_token} signed RS256; Spring's OAuth2 client verifies it
against Google's JWKS at
\code{https://www.googleapis.com/oauth2/v3/certs}. Google never hands you
its private key, and you never need it.

\subsection{The costs}

\begin{center}
\begin{tabular}{@{}lll@{}}
\toprule
 & HS256 & RS256 (2048-bit) \\
\midrule
Key material & one secret, $\ge$ 32 bytes & private + public, $\approx$1.7\,KB PEM each \\
Signature size & 32 bytes (43 chars) & 256 bytes (342 chars) \\
Sign & microseconds & about 1\,ms \\
Verify & microseconds & tens of microseconds \\
Who can sign & any key holder & private-key holder only \\
Who can verify & any key holder & anyone \\
\bottomrule
\end{tabular}
\end{center}

RSA verification is far cheaper than signing because $e = 65537$ is tiny
while $d$ is 2048 bits wide. That asymmetry fits the topology: one signer,
many verifiers. If signature size matters, ES256 (ECDSA on curve P-256)
gives 64-byte signatures and much smaller keys, and is the modern
default for new identity providers.

\begin{pitfall}[algorithm confusion]
Symptom: a token that should be rejected is accepted --- and your logs
show nothing wrong. Cause: older libraries read the \code{alg} field of
the header to decide \emph{how} to verify. An attacker takes a legitimate
RS256 token, rewrites the header to \code{"alg":"HS256"}, and signs it
with HMAC using the \emph{public key bytes} as the secret --- bytes that
are, by design, public. A naive verifier sees HS256, uses the public key
it has as an HMAC secret, and the check passes. JJWT prevents this by
letting the \emph{key type} choose the algorithm, not the header:
\code{verifyWith(SecretKey)} rejects any asymmetric token and
\code{verifyWith(PublicKey)} rejects any HMAC token. That is why
\code{JwtAuthenticationFilter} is safe today, and why you must never
verify with a method that accepts ``whatever the header says''.
\end{pitfall}

\subsection{Converting the project}

The change is small enough to read in full. Generate the pair once:

\begin{lstlisting}[language=bash]
$ openssl genpkey -algorithm RSA -pkeyopt rsa_keygen_bits:2048 \
      -out jwt-private.pem
$ openssl pkey -in jwt-private.pem -pubout -out jwt-public.pem
\end{lstlisting}

The private PEM goes only into auth-service's Secret; the public PEM can
go into a ConfigMap, because it is not secret. \code{JwtService} signs
with the private key:

\begin{lstlisting}[language=Java]
@Service
public class JwtService {

    private final PrivateKey privateKey;   // sign
    private final PublicKey publicKey;     // self-verify where needed

    public JwtService(@Value("${app.jwt.private-key-pem}") String privPem,
                      @Value("${app.jwt.public-key-pem}") String pubPem)
            throws GeneralSecurityException {
        KeyFactory kf = KeyFactory.getInstance("RSA");
        this.privateKey = kf.generatePrivate(
                new PKCS8EncodedKeySpec(decodePem(privPem)));      (*@\co{1}@*)
        this.publicKey = kf.generatePublic(
                new X509EncodedKeySpec(decodePem(pubPem)));
    }

    public String generateAccessToken(User user) {
        return Jwts.builder()
                .header().keyId("2026-09-key1").and()             (*@\co{2}@*)
                .subject(String.valueOf(user.getId()))
                .claim("email", user.getEmail())
                .claim("role", user.getRole().name())
                .issuedAt(new Date())
                .expiration(new Date(System.currentTimeMillis()
                        + accessTokenExpiryMs))
                .signWith(privateKey, Jwts.SIG.RS256)             (*@\co{3}@*)
                .compact();
    }

    private static byte[] decodePem(String pem) {
        String body = pem.replaceAll("-----BEGIN [A-Z ]+-----", "")
                         .replaceAll("-----END [A-Z ]+-----", "")
                         .replaceAll("\\s", "");
        return Base64.getDecoder().decode(body);
    }
}
\end{lstlisting}

\begin{description}
  \item[\co{1}] PKCS\#8 is the container format \code{openssl genpkey}
  writes for private keys; X.509 \code{SubjectPublicKeyInfo} is what
  \code{-pubout} writes. Strip the PEM armor, base64-decode, done.
  \item[\co{2}] \code{kid} (key id) names the key that signed this token.
  It costs a few bytes now and makes rotation possible later.
  \item[\co{3}] JJWT infers RS256 from the key type; naming it
  explicitly documents the intent.
\end{description}

The gateway --- and any service that wants to verify for itself ---
holds only the public key:

\begin{lstlisting}[language=Java]
public JwtAuthenticationFilter(@Value("${app.jwt.public-key-pem}") String pubPem)
        throws GeneralSecurityException {
    this.publicKey = KeyFactory.getInstance("RSA")
            .generatePublic(new X509EncodedKeySpec(decodePem(pubPem)));
}

Claims claims = Jwts.parser()
        .verifyWith(publicKey)      // PublicKey: accepts RS/PS/ES only
        .build()
        .parseSignedClaims(token)
        .getPayload();
\end{lstlisting}

Nothing else changes: same claims, same headers, same
\code{GatewayHeaderAuthFilter}. The gateway can no longer mint tokens,
and \code{api-gateway.yml} no longer needs \code{app.jwt.secret} at all.

\subsection{Rotation, and what \code{kid} is for}

Keys age: policy says rotate every 90 days, or a laptop with the PEM on it
is lost. Rotation without logging every user out is a four-step ceremony:

\begin{enumerate}
  \item Generate a new pair; add its public key to the JWKS alongside the
  old one, under a new \code{kid}.
  \item Wait until every verifier has refreshed its cache (or restart
  them).
  \item Switch auth-service to sign with the new private key. New tokens
  carry the new \code{kid}; old tokens, still in flight, carry the old
  one --- and verifiers pick the matching public key by \code{kid}.
  \item After the longest token lifetime has elapsed (here 15 minutes
  for access tokens), remove the old key from the JWKS.
\end{enumerate}

Without \code{kid} a verifier would have to try every key it knows;
with it, the lookup is exact and the old key can be retired with
confidence. The project's static-PEM version above has no JWKS endpoint;
the natural next step is to serve \code{/.well-known/jwks.json} from
auth-service and have verifiers use Spring Security's
\code{NimbusJwtDecoder.withJwkSetUri(...)}, which caches and refetches on
an unknown \code{kid} automatically.

\section{Payload bloat}
\label{sec:payload-bloat}

Statelessness is seductive: if the token carries everything a service
needs, no service ever calls another to ask ``may this user do this?''
Teams follow that logic until the token carries the user's roles,
permissions, groups, feature flags, and profile --- and weighs 4\,KB. The
arithmetic then turns against them. A dashboard page firing 20 parallel
API calls ships 80\,KB of token before a single byte of business data.
Multiply by tens of thousands of concurrent users and the authentication
header is a measurable share of ingress bandwidth. Cookies, if that is
where the token lives, have a 4\,KB ceiling; some proxies cap header size
at 8\,KB; logs that record headers bloat in step.

The project's access token is thin. \code{JwtService.generateAccessToken}
writes \code{sub}, \code{email}, \code{role}, \code{iat}, \code{exp} ---
about 200 bytes signed. Keep it that way by rule: \textbf{a token carries
identity; entitlements are looked up.} The moment course-service needs to
know ``which courses may this teacher edit?'', the temptation to add a
\code{courseIds} claim appears. Resist it. That list changes when a
course is created, which the token cannot know until it expires; and it
grows without bound. The service already owns the enrollment data
(Chapter~7) --- one indexed query answers the question correctly and
currently.

\section{The phantom token pattern}

Suppose the fat token is unavoidable --- a legacy authorization model, a
partner integration. The phantom token pattern keeps the \emph{external}
token tiny while the \emph{internal} one stays self-contained:

\begin{enumerate}
  \item Auth-service creates the full JWT but does \emph{not} return it.
  It stores the JWT in a fast store (Redis) under a random opaque
  identifier, and returns only that identifier --- a UUID, 36 bytes.
  \item The gateway receives the UUID, looks it up, and swaps in the real
  JWT before forwarding. Internal services see a normal signed token.
  \item Internal services verify it themselves (RS256, JWKS). Inside the
  perimeter the system is stateless and zero-trust; outside, the client
  carries a token with no readable content at all.
\end{enumerate}

Logout becomes trivial and \emph{immediate}: \code{DEL <uuid>} in Redis,
and the next request with that identifier fails at the gateway with 401,
no matter how long the inner JWT would have lived. This is the pattern's
real prize --- it repairs JWT's one structural weakness, irrevocability.

The project already runs half of it. Look at
\code{JwtService.generateRefreshTokenValue}:

\begin{lstlisting}[language=Java]
public String generateRefreshTokenValue() {
    return UUID.randomUUID().toString();
}
\end{lstlisting}

The refresh token is not a JWT. It is an opaque identifier stored in the
\code{refresh\_tokens} table (Flyway \code{V2}) and revoked by deleting
the row --- \code{AuthService.logout} does exactly that. The access token
remains a self-contained JWT. So the project's revocation latency is
bounded by the access token's 15 minutes, and its revocation
\emph{mechanism} is the phantom pattern applied to the long-lived
credential only.

\begin{pitfall}[refresh without rotation]
Read \code{AuthService.refreshToken}: it returns the \emph{same} refresh
token it was given. A refresh token stolen once is therefore valid for
its full seven days, and the theft is undetectable. The industry
practice is \emph{rotation}: every \code{/refresh} issues a new refresh
token and invalidates the old one. If an old token is ever presented
again, two parties hold copies --- revoke the whole family. The table
already has the columns to do this; the change is a delete-and-insert in
one transaction.
\end{pitfall}

\begin{decision}[three ways to revoke an access token]
When an admin disables an account, the project's answer is ``within 15
minutes''. The alternatives, in ascending cost:
\begin{enumerate}
  \item \textbf{Wait it out.} Short access tokens plus stateful refresh
  tokens. Zero extra infrastructure; revocation lag equals access-token
  lifetime. Acceptable for most applications, and what the project does.
  \item \textbf{Deny-list.} Add a \code{jti} (token id) claim; on
  logout or disable, write the \code{jti} to Redis with a TTL equal to
  the token's remaining life; the gateway checks Redis on every request.
  One fast lookup per request, and only revoked tokens are stored.
  \item \textbf{Phantom token.} Opaque identifier outside, JWT inside,
  Redis as the mapping. Immediate revocation and no readable token in
  the browser --- but Redis is now in the critical path of \emph{every}
  request. Redis down means every user is logged out at once.
\end{enumerate}
A senior engineer's objection to the article's enthusiasm: option~3
trades ``cannot revoke'' for ``one new single point of failure''. Most
teams should reach for option~2 first and treat option~3 as a fit for
banking-grade requirements, where the article itself admits sessions may
be the better tool anyway.
\end{decision}

\section{What breaks at 3\,a.m.}

Failure analysis is the fastest way to tell the two key models apart.
Auth-service goes down at 3\,a.m.; nothing else does.

\begin{itemize}
  \item \textbf{HS256, as deployed.} The gateway holds the secret and
  verifies locally: every logged-in user keeps working until their
  access token expires. Refresh calls fail (auth-service is the
  \code{/refresh} endpoint), so sessions die within 15 minutes. New
  logins fail. Recovery is instant when auth-service returns.
  \item \textbf{RS256 with cached JWKS.} Identical for running
  verifiers: they cached the public key at startup and need nothing from
  auth-service. The new failure is a \emph{cold} verifier --- a gateway
  pod restarted during the outage cannot fetch the JWKS, and depending on
  its configuration either refuses all tokens (fail closed) or crashes
  at boot. Mitigations: ship the current public key in configuration as
  a fallback, or give the JWKS cache a long grace period.
  \item \textbf{Phantom token.} Add Redis to the list of things that must
  be up. Auth-service down has the same effects as above; Redis down is
  a total outage of authenticated traffic.
\end{itemize}

The lesson generalizes: statelessness buys you independence from the
issuer at request time. Every mechanism you add to regain control ---
JWKS, deny-lists, opaque mappings --- reintroduces a dependency, and you
should be able to name it and say what happens when it is unavailable.

\section{When JWT is the wrong tool}

The article closes with a warning worth repeating, because the project's
own path (Chapter~27) includes a monolithic alternative. JWT earns its
complexity when \emph{many independent parties} must answer ``who is
calling?'' without sharing a session store: microservices, single
sign-on across applications, OAuth2/OIDC federation. It is the wrong
tool when the system is a single deployable with room for its sessions in
memory, when permissions change often and must take effect immediately,
or when a regulator requires that a session be terminable in
milliseconds. In those cases a server-side session with a cookie is
simpler, revocable, and boring --- and boring is a feature.

\section{Outside the project}

Keycloak, Auth0, Cognito, Okta and Entra ID all issue RS256 or ES256
tokens and publish JWKS; every Spring resource server in the industry
verifies with the one-line \code{jwk-set-uri} configuration shown in
Chapter~5. Service meshes (Istio, Linkerd) implement the mTLS variant of
zero trust so that application code keeps reading headers while the
transport proves who wrote them. The phantom token pattern is
productized by Curity and by API gateways such as Kong and Apigee under
names like ``token introspection'' and ``opaque token exchange''. Knowing
the hand-rolled version --- one HS256 secret, two holders --- is what
makes those products legible.

\section{Questions from real interviews}

These are asked as open problems, not definitions. The answers below are
what a senior engineer says in the first two minutes; the interviewer
then digs into whichever trade-off you named.

\subsection*{Q: How do you invalidate a JWT before it expires?}

You do not; you bound the damage. Access tokens live minutes, not hours,
so ``wait it out'' is the baseline. For immediate revocation add a
\code{jti} claim and a deny-list in Redis with TTL equal to the token's
remaining life --- one lookup per request, only revoked tokens stored,
and Redis unavailability can fail open or closed by policy. Full
opacity (phantom token) makes revocation instant but puts Redis on
every request's critical path. In this project the refresh token is
already opaque and database-backed, so ``log out everywhere'' is
\code{DELETE FROM refresh\_tokens WHERE user\_id = ?} plus, at most, 15
minutes of residual access.

\subsection*{Q: Login traffic goes from 10 to 10\,000 per second. What breaks?}

Signing first: RS256 signing is about 1\,ms of CPU per token, so ten
thousand logins per second is ten cores of pure signing, while HS256
would be negligible. Then password hashing, which is deliberately slow
(bcrypt at cost 10 is roughly 100\,ms) --- that is the real wall, and
the answer is horizontal scaling of auth-service plus rate limiting
per IP and per account at the gateway. Verification is not the problem:
RS256 verify is tens of microseconds and needs no network. Also check
that the refresh-token insert is not on a single hot index and that
the \code{refresh\_tokens} table is pruned of expired rows.

\subsection*{Q: Two regions, one auth system. What has to be shared?}

For HS256, the secret, which is a distribution problem you do not want
across regions. For RS256, only the public key set, which is
cacheable, public, and can lag safely for the length of the rotation
overlap. Clocks matter more than keys: \code{exp} and \code{iat} are
compared to the verifier's clock, so allow a small skew (JJWT's
\code{clockSkewSeconds}, typically 30--60\,s) and run NTP everywhere.
Refresh tokens are stateful, so either pin refresh to the issuing
region or replicate the table; the project's single Postgres would be
the first thing to change.

\subsection*{Q: A token was stolen from a user's browser. What do you do?}

Contain, then detect. Contain: revoke the user's refresh tokens; the
access token dies within its lifetime. Detect: refresh-token rotation
turns a replayed old token into a signal --- two parties presenting the
same family means theft, revoke the family. Prevent: keep the refresh
token in an \code{HttpOnly} cookie so injected scripts cannot read it,
and keep the access token in memory rather than \code{localStorage}
(Chapter~38). Add device or IP binding only if you accept the
false-positive cost for mobile users.

\subsection*{Q: Why not just check the database on every request?}

You can, and monoliths do --- that is a session. The reason not to in
a service architecture is that every service would need the user store
or a call to the one that has it, which puts auth-service on the
critical path of everything and couples every deploy to it. A signed
token lets each hop answer ``who is this'' locally. The honest answer
names the cost --- revocation lag --- and the price you paid to buy
back control: short lifetimes, a stateful refresh token, optionally a
deny-list.

\subsection*{Q: What is in your token, and what is deliberately not?}

Subject, role, issued-at, expiry, and here an e-mail --- identity, not
entitlements. Not in it: permissions lists, group memberships, anything
that changes more often than the token expires, anything another
service owns, and anything secret (the payload is readable). If an
interviewer pushes for ``put the permissions in so we skip a query'',
the senior answer is the arithmetic of Section~\ref{sec:payload-bloat}:
4\,KB times every request, and a stale authorization decision
frozen for the token's lifetime.

\begin{quiz}
\begin{enumerate}
  \item After the RS256 conversion, \code{api-gateway.yml} no longer
  contains \code{app.jwt.secret}. If the entire config-server were
  leaked, could the attacker forge a token? Explain which file would
  have to leak instead.
  \item Describe the 90-day rotation ceremony and state precisely which
  problem \code{kid} solves in step~3.
  \item Thirty services each verify RS256 tokens with a JWKS fetched at
  startup. Auth-service dies at 3\,a.m. Which users are affected, when,
  and which single event during the outage would make a verifier reject
  \emph{all} traffic?
  \item A colleague proposes adding a \code{courseIds} claim so that
  course-service can authorize without a query. Give two reasons to
  refuse that are independent of token size.
  \item The gateway forwards a request to a public path without touching
  headers. Explain why this is a latent rather than an active
  vulnerability today, and name the one-line configuration that removes
  the latency.
  \item An admin disables a user who is logged in. Give the revocation
  lag under the project's current design, and name the extra dependency
  each of the two faster designs introduces.
\end{enumerate}
\end{quiz}
